\input{./._relpath-to-latexroot.ltx}
\documentclass[\latexroot/\projectname]{subfiles}
\input{\latexroot/@resources/subfile-setup.ltx}
\begin{document}

\section{Inspecting the mechanism}
\whenintegrated{\label{sec:mechanism}}

In this section we perform a series of simpler experiments to understand the forces that underlie the positive and welfare results for the sudden stop. We consider three scenarios: (i)~a real benchmark with home-good denominated debt, (ii)~flexible exchange rates with foreign-currency denominated debt, and (iii)~fixed exchange rates with foreign-currency denominated debt.

The key finding is that the distribution of household leverage is a crucial determinant of aggregate dynamics. When poor households are more highly leveraged, a larger contraction in aggregate consumption takes place, because the fall in net wealth hits households with high marginal propensities to consume. The desirability of a fixed versus floating exchange rate depends crucially on the distribution of household leverage and the fraction of debt denominated in foreign currency.

% Smart bibliography: Only include bibliography if standalone AND has citations
\smartbib

\end{document}
